\textbf{Задача 8} Докажите, что $\np\neq\e$.

\begin{proof}

\begin{enumerate}
    \item Покажем, что $\e\subsetneq\EXP$. Вложение следует напрямую из определения классов.
    
    Теперь воспользуемся теоремой:

    \textbf{Теорема} (об иерархии по времени). $t(n)$ называется конструируемой по времени, если можно вычислить $t(n)$ на входе $1^n$ за время $O(t(n))$. Если $f$ и $g$ строго монотонные конструируемые по времени функциив, при этом $f(n)\log f(n)= o(g(n))$, то $\dtime{f(n)}\neq\dtime{g(n)}$. 
    
    Покажем строгость вложения, то есть неравенство классов:
    
    \[\forall c \quad 2^{cn} \cdot \log\left(2^{cn}\right) = cn2^{nc} = o\left(2^{n^2}\right)\]
    
    значит, $\e \subsetneq \dtime{2^{n^2}} \subset \EXP$
    
    \item Предположим, что $\e = \np$. 

    Пусть произвольный $A \in \dtime{2^{n^c}} \subset \EXP$.
    
    Объявим класс $A_{pad} = \left\{\left(x,1^{|x|^c}\right), x \in A \right\}$.
    Новый $x' = \left(x,1^{|x|^c}\right)$ имеет длину $n' = n + n^c = O(n^c)$.
    
    Значит, $A_{pad} \in \dtime{2^{n'}}$, то есть $A_{pad} \in \dtime{2^{n}} \subset \e = \np$. Значит, $A_{pad} \in \np$.
    
    Но тогда $A \in \np$, так как ко входу $x \in A$ можно приписать $1^{|x|^c}$ и запустить машину для $A_{pad}$.
    
    В силу произвольности $A$, получаем $\np=\EXP$, но это противоречит тому, что $\e \neq \EXP$. Значит, $\e \neq \np$. 
\end{enumerate}

\end{proof}
